\section{User Interfaces for Information Access}
\label{sec:iaui}
\subsection{Contribution: DAVE}

\subsection{entity-centric architecture}
from claw. also part of wise
\todo{dave and its evolution through the phd and papers}


\subsection{intro}

In this paper, we introduce DAVE, a tool for assisted analysis of document collections in knowledge-intensive domains. 
The tool goal is to support search needs that span across different points of the extractive vs. abstractive spectrum, as discussed in~\cite{worledge2024extractive}.     
It features a graphical user interface (GUI) that enables users to visualize, explore, and query documents.
% using the established paradigms of faceted search and conversational interfaces.
%DAVE stands out for several key reasons. It provides all these features within a single tool: %(i) it focuses on the exploration of documents and annotations rather than Information Extraction itself, ensuring flexibility and adaptability in the rapidly evolving AI landscape; 
Our tool is specifically tailored  for domains where entities are first-class citizens in document analysis and provide the main following features:
(i) an entity-driven faceted search interface for entity-driven exploration; 
(ii) a conversational interface supporting complex natural language queries; 
(iii) seamless integration of faceted search and conversational interaction to refine document sources;
(iv) a human-in-the-loop mechanism for refining entity annotations, ensuring corrections are propagated throughout the system.
%, enhancing the analysis of large and complex document collections; 
%(v) it enables search and exploration at different levels of granularity: individual documents, entire corpora, and sub-corpora; 

%DAVE has already been tested \cite{pozzi_wise_2025,BELLANDI2024105904} in real-world use cases within the legal domain, receiving positive feedback from domain experts and end-users for its usability and exploration capabilities.
DAVE is designed as an open system and is released as open source under the Apache-2.0 license. Documentation, source code, and demonstration videos are available on the project GitHub page\footnote{\url{https://github.com/unimib-datAI/DAVE}}.

\subsection{related}
DAVE has been used in prototypes for Italian projects in the legal domain~\cite{BatiniSPBFZPVR24} with the goal of showing stakeholders AI-powered search functionalities, we have experimented DAVE in i) search on court decision in civil trials~\cite{BELLANDI2024105904}, ii) criminal investigations and chat analysis~\cite{pozzi_wise_2025}, and iii) analysis of the documentation about the Bologna massacre of August 2, 1980.


% Two of these scenarios contain sensitive data in Italian, which stressed the need to develop a solution with on-premise small LLMs.   

In relation with related work, we discuss below the three main novelties that we believe DAVE presents as a system: 

%\begin{itemize}
    %\item Mixing entity-driven faceted search and chatbot for information access; Q: is there any work on combining those two access mechanisms? (H: entity-driven faceted search is somehow mainstream technology in many knowledge-driven scenarios where documents are analyzed using entity extraction (EE) methods

    %\item Integrating interactive entity-driven knowledge consolidation in an information exploration interface; Q: any interactive approach to entity co-referencing? H: several applications have been proposed to support text annotations, e.g., Doccano~\cite{doccano} is one of the most used platforms; very recent work on how to improve annotations using interactive approaches and minimizing users' effort; while we have not fully integrated sophisticated processes to improve the annotation quality as proposed in~\cite{klie-etal-2020-zero}, our application is the first one to incorporate interactive editing methods into an \textit{explorative search} interface; our work is fully in agreement with recent approaches that propose resource-efficient interactive methods to improving annotations; in fact, inspired by a pay-as-you-go paradigm~\cite{klie-etal-2018-inception}, we use annotations computed by automatic methods but provide users consuming processed documents with functionalities to correct and improve algorithmic results. While we did not integrate sophisticated mechanisms to refine NER and NEL results yet, we instead developed a method to refine entity clustering, one of the most challenging task in end-to-end EE pipelines.      

%\end{itemize}

\begin{itemize}
    \item \textit{Mixing entity-driven faceted search and conversational assistant}.
    Entity-driven faceted search is a mainstream technology in many knowledge-driven scenarios where documents are analyzed using entity extraction (EE) methods~\cite{grafs,hirsch-etal-2021-ifacetsum}.
    On the other hand, RAG systems on pre-filtered data have been studied, with approaches like agent-based filtering ~\cite{poliakov2024multi}, metadata-based filtering ~\cite{chang2024main}, and natural language inference ~\cite{yoran2023making}. Our work combines faceted search with a RAG system, enabling dynamic, entity-driven document filtering.
    
    \item \textit{Integrating interactive entity-driven knowledge consolidation in an information exploration interface}.
    Several platforms support text annotations, with Doccano~\cite{doccano} being widely used; 
    Very recent work has surveyed interactive approaches to improve annotations, minimize user effort, manage annotation teams, support pre-annotated data, and enable customizable task design~\cite{de2024integrating}.
    While we haven't fully integrated the advanced annotation quality methods from ~\cite{klie-etal-2020-zero}, our application is the first to integrate incorporate interactive methods editing into an \textit{exploratory search} interface, focusing on improving entity clustering, a key challenge in end-to-end EE pipelines.


%nesusn approccio exploration search integra editing
% e nessun aprpoccio di editing integra esploration search
%--> si vede che no ha senso... e invece si, è il nsotro use case 
%permette ai giudici di esplorare i dati annotati da una pipeline e di avere comquneu potere di correggere per migliorare la qualità del tool di exploration

%Although semantic annotation tools has a rich list of proposals addressing several aspects of the process as surveyed in ~\cite{de2024integrating}  such as management of annotation teams, support for pre-annotated data, customizable task design among others, no approch have integrated interact drive knowledge consolidation in a information exploration interface

    \item \textit{Entity-centric RAG prototyping and grounding}.
    Frameworks like LangChain ~\cite{langchain} facilitate rapid prototyping and configuration of RAG systems, while tools such as RAGAS ~\cite{es-etal-2024-ragas} and RAGChecker ~\cite{ru2024ragchecker} allow for detailed evaluation through a wide range of metrics.
    However, existing applications do not support the prototyping of highly entity-centric LLM-RAG systems that enable direct analysis of the corpus and its entities to verify factual accuracy, making effective debugging more challenging.
\end{itemize}

%    \item Data linking for knowledge management 

These innovations arise from the need for domain experts to thoroughly analyze and explore documents annotated by entity extraction pipelines, using established, user-friendly search paradigms and allowing experts to improve the system by refining annotations. A first quantitative evaluation where DAVE outputs are compared to outputs of top-tier models is ongoing in the context of a Civil Appeal Proceedings use case~\cite{ldql}.   

During the demonstration session, users are guided in exploring a document collection using DAVE's features.


\subsection{framework}


%DAVE approach to document analysis is centered on entities. Entities in DAVE are leveraged in two main ways. the first one is as filter for facedeted search, the second one is visualization of document interface

% Knowledge intensive domains
Knowledge-intensive domains feature extensive collections of complex documents where information is highly factual and deeply connected to real-world entities. Moreover, details about a single entity are often distributed across multiple documents within the corpus. 
%As a result, there are significant challenges in efficiently addressing sophisticated questions, especially when relying on manual methods or tools that are not specifically tailored for entity-driven analysis.
%The role of entity annotations in DAVE
Recognizing that entities are the cornerstone of effective document analysis in knowledge-intensive domains, DAVE proposes an entity-centric exploration approach that %. Our entity-centric approach
goes beyond the exploitation of entity mentions, e.g., as an output of a NER algorithm, supporting mentions linked to entity identifiers and clusters of mentions referring to the same entity.
% As a result, a mention is linked to an entity, and an entity is linked to possibly more mentions.
% This approach, while compatible with simpler IE extraction approaches, is thought to be aligned with end-to-end entity extraction pipelines.
% that account for entity clustering in a direct (e.g., via coreference resolution~\cite{logan2021benchmarking}) or indirect (e.g., via entity linking~\cite{pozzi_wise_2025,kassner-etal-2022-edin}) way.      

  
DAVE provides users with a number of functionalities to address the outlined requirements, which are discussed with examples in legal documents exploration (Figure \ref{fig:dave}):
\begin{itemize}
    \item \textbf{Search \& Filter}: Users can retrieve documents by entering keywords, which is ideal when they have a clear idea of what they are looking for. Additionally, they can refine their search results by applying multiple filters. For instance, users can filter the corpus based on specific individuals and locations.
    Figure illustrates the \textit{Search \& Filter} functionalities (the blue section), which allow users to retrieve documents by keyword and refine results through entity-based filters. For example, a user reviewing documents from a trial might search for the keyword ``appeal''. To narrow the focus, they could use the filtering panel on the left, which organizes filters by entity type (e.g., person, location). By selecting ``Licio Gelli'' from the list of identified entities, the user can dynamically refine the results, isolating only those documents where his name appears, with both the result set and available filters updating accordingly.
    
    \item \textbf{Explore}: Users can browse the full list of entities identified in each document and navigate directly to sections where they are referenced.
    In the figure, after selecting a document in the \textit{Search \& Filter} interface, it can be explored in detail via the \textit{Explore} interface (green section). The entire document is shown with highlighted entity mentions based on type, and a list of entities grouped by mention appears on the left. 
    For example, in the figure, the entity ``Valerio Fioravanti'' is mentioned 4 times. Clicking any mention takes the user directly to the corresponding text span.
    
    \item \textbf{Conversational Question Answering (QA)}: 
    Natural language queries are supported, allowing users to ask questions about entities and factual information across multiple documents. In the example (yellow section), the user asks the chatbot about the appeal grounds raised by Prosecutor Squadrini in the Musumeci and Belmonte case, and the DAVE chatbot responds with relevant document passages.
     
    \item \textbf{Knowledge Refinement}: Users can refine entity clusters, ensuring that corrections are reflected across all system functionalities. In the figure (red section), the user has identified that the entities ``Corte d'Assise di Bologna'' and ``Corte di Assise di Bologna'' with their respective mentions are actually equivalent, therefore is collapsing the two entity clusters into one.
\end{itemize}

DAVE also supports \textit{search composability}, enabling users to combine Search, Filter, and Conversational Question Answering. By filtering results before querying the RAG-LLM system, the chatbot works with a more focused document set, improving answer precision and relevance. Additionally, DAVE ensures data privacy through on-premise servers and access controls, crucial for sensitive domains.% like legal and healthcare.

By offering this suite of functionalities, DAVE enables precise control and in-depth understanding of large corpora, providing diverse exploration and search capabilities tailored to knowledge-intensive domains and specialized users. To this aim we implemented these functionalities by mixing the following techniques and technologies: 

\begin{itemize}
    \item \textbf{Entity-centric management}. Entity-level annotations are represented in the GATE format~\cite{gate2002} and stored in a MongoDB database; the DB stores information about the annotations for every entity mentions, entity identifiers, and links between annotations and entity identifiers. The current prototype considers annotations resulting from pipelines that apply algorithms for NER, Named Entity Linking (with links to Wikipedia), NIL Prediction (to identify entities not represented in Wikipedia), and NIL Clustering (to cluster NIL entities and create identifiers for each cluster)\cite{pozzi2022,PozziRBP23,BELLANDI2024105904}. As a result, each entity mention can be linked to an entity identifier, either external (e.g. a Wikipedia URI), or internal (identifying a local cluster). Each cluster is associated with a default surface form used to display the entities in the interface.  
    
    % \item \textbf{Search \& Filter}: The \textit{Search} functionality is powered by a search engine, which uses keyword matching to efficiently retrieve relevant documents. This search engine serves as the foundation for several other functionalities, providing fast and efficient document retrieval.
    \item \textbf{Keyword and Faceted Search}: The \textit{Search} engine uses keyword matching to efficiently retrieve relevant documents and serves as the foundation for several other functionalities.
    %\item \textbf{Faceted Search}: 
    To support the \textit{Search \& Filter} and \textit{Explore} functionalities, DAVE employs the well-established faceted search paradigm. This technique provides filtering facets based on entities, allowing users to refine their search results. By providing a structured way to narrow down results, faceted search enhances the user's ability to explore large corpora. Users can apply filters either across the entire corpus or within a subset of documents for a more granular exploration.
    
    \item \textbf{Human-in-the-Loop (HITL)}: \textit{The Knowledge Consolidation} feature follows the HITL paradigm, ensuring continuous user involvement in refining the system. Users can correct and refine annotations and entity clusters, and these corrections are reflected across all system functionalities. This active participation helps the system improve over time, ensuring a pay-as-you-go consolidation of the background data as proposed for similar tasks~\cite{de2024integrating,de2024interactive,cutrona2019asia,cruz2016quality}.
    %used by background knowledge base data as that the generated results are aligned with user expectations.
    \item \textbf{Retrieval Augmented generation (RAG)}: The Conversational QA functionality is powered by an LLM-based chatbot, implemented through the RAG paradigm. This enables users to ask fact-based queries about entities across multiple documents. While LLMs excel in natural language understanding and zero-shot learning, RAG ensures responses are grounded in retrieved documents, addressing concerns about hallucinations and limited knowledge, which is crucial in certain domains.
\end{itemize}



%The most basic search fucntionality (i.e., key-word search) is supported by a search engine. Beyond this simple use of an index text, this backen search-engine feed several other functionalities.
%
%To support the \textit{filter} and \textit{explore} functionalities, DAVE employs the well-established faceted search paradigm, organizing facets by entity types. This approach enables users to filter the corpus while gaining an overview of the entities it contains, facilitating an initial exploration phase. Faceted search can be applied to a subset of documents or a single document. In the latter case, facets offer highly specific insights, enhancing document-level exploration. Document filtering to work on semantic annotations provided as input with the corpus and on the search engine to filter data by entities. 
%
%Concerning the \textit{knowledge consolidation} feature it follows the Human-In-The-Loop (HITL) paradigm ensuring active human involvement in the process, directly contributing to the overall improvement of the system's performance as corrected annotations reverse in the other functionalities. 
%The \textit{ask question} feature is implemented as a conversational interface powered by Retrieval-Augmented Generation (RAG) using Large Language Models (LLMs). This interface allows users to formulate complex queries in natural language, with the interpretation of query intent handled by the LLM. However, despite their exceptional capabilities in text comprehension and zero-shot learning, LLMs present two major challenges in knowledge-intensive domains. First, their knowledge is limited to the data they were trained on, making them unable to access or reason over external or updated information. Second, standard LLMs are not designed for search tasks, as they can generate hallucinated responses. In domains where factual accuracy is critical, these limitations can significantly impact reliability. By combining LLMs with the RAG paradigm, DAVE mitigates these risks, grounding responses in retrieved documents while providing a high-performance conversational assistant. 
%
%Finally, DAVE’s key strength lies in its ability to access extracted knowledge through a comprehensive faceted search, integrated with a conversational search feature. This conversational search can be applied to the entire corpus or used in conjunction with the semantic search results, utilizing them as a filtering mechanism. This functionality is enabled by the integration of the RAG-LLM system with a search engine that preserves the relationships between document chunks, full documents, and entities, thus allowing for precise filtering of the RAG-LLM search scope.



%is backed by semantic annotations
%- semantic annotations
%- keyword search
%- faceted search sfruttando annotazioni semantiche
%- chatbot, RAG

%core features
%\begin{itemize}
%    \item access control: imporante per domini che trattano dati sensibili
%    \item semantic exploration
%    \item human editing (annotations, clustering)
%    \item Free Q/A via Chatbot

%Why DAVE:
%è uno strumento E2E che mette insieme la potenza di pipelien IE, l'intevernto umano e la consultazione intelligente della conoscenza attraverso
%e laa potenza dei LLM focalizzati sui dati caricati e filtrati sulal ricerca semantica
    %attraverso una potente interfaccia utente permette di consultare la conoscenza estratta


%perchè DAVE spacca
%The strength of DAVE lies in its ability to consult extracted knowledge through a full-featured semantic search, combined with a conversational search capability. This conversational search can operate across the entire corpus or downstream of the semantic search, leveraging the latter as a filtering mechanism.

% perchè abbiamo scelto questo tipo di interfacce
%The semantic search, implemented using the well-known faceted search mechanism, provides an intuitive filter system, widely adopted in domain-specific search engines such as those used in e-commerce, thanks to its effectiveness and ease of use. Similarly, the conversational search, built on the RAG paradigm, offers an exceptionally user-friendly and flexible interface. It also ensures source attribution in its responses, effectively mitigating the phenomenon of hallucinations.

%perchè queste scelte sui paradigmi di ricerca sono rilevanti nel legalAI
%In the legal domain, the combination of advanced search capabilities and ease of use is particularly critical. Legal professionals often work under time constraints, needing to quickly retrieve precise and reliable information from vast amounts of complex text. 
%Semantic search facilitates this by enabling intuitive filtering and targeted exploration,  while conversational search offers a natural, interactive way to query and retrieve information. 
%These features reduce time spent searching and ensure transparency through source attribution, addressing the critical need for efficiency, accountability, and reliability in the legal domain.



